\chapter{Tái lập thí nghiệm}
\label{app:reproduce}

\section{Mã và dữ liệu sử dụng}
Mã BFS trong \href{https://github.com/duyanhtr130905/parallel-bfs-openmp}{kho mã Parallel BFS}
được lấy tại commit
\href{https://github.com/duyanhtr130905/parallel-bfs-openmp/commit/fb990bf2ddb3adb6bac5e218ef32fd40adfc35cd}{\texttt{fb990bf}}.
Lịch sử dự án lưu bộ dữ liệu thực nghiệm khác để đối chiếu khi cần.
\begin{itemize}
\item Trình đo \path{scripts/benchmark_six_graphs.c}.
\item Số liệu \path{results/experiment-20261124/raw.csv}.
\item Trình đo bổ sung \path{scripts/benchmark_supplement.c} và trình sinh đầu vào \path{scripts/run_supplement.py}.
\item Kế hoạch, manifest, CSV theo đồ thị và hồ sơ frontier trong \path{results/experiment-20261125/}.
\end{itemize}

\section{Lệnh biên dịch và chạy trên máy của nhóm}
Đặt trình đo cùng cây mã BFS nêu trên và chạy từ thư mục gốc. Trên Windows, dùng GCC và libgomp của MSYS2 UCRT64. Lệnh PowerShell sau ghi tệp mới cho phiên lặp.
\begin{quote}\small\singlespacing
\begin{verbatim}
$flags = '-O2','-std=c11','-Wall','-Wextra','-Werror',
    '-fopenmp','-Iinclude','-Isrc'
$files = 'scripts/benchmark_six_graphs.c','src/graph.c',
    'src/bfs_serial.c','src/bfs_parallel_basic.c',
    'src/bfs_parallel_optimized.c'
gcc @flags @files -o benchmark_six_graphs.exe
if ($LASTEXITCODE -ne 0) { throw 'Loi bien dich' }
cmd /c ".\benchmark_six_graphs.exe > benchmark_repeat.csv"
if ($LASTEXITCODE -ne 0) { throw 'Loi phep do' }
\end{verbatim}
\end{quote}
Khi lặp phép đo, ghi lại cấu hình thực tế tại thời điểm chạy và giữ từng tệp CSV theo phiên riêng. Trên một máy khác, kết quả thu được thuộc một môi trường khác và cần trình bày riêng.

Phiên 25/11 dùng mã BFS tại \texttt{0d5a53b}. Để lặp phiên này, chạy
\path{scripts/run_supplement.py} với Python, \texttt{psutil} và GCC có OpenMP.
Tham số \texttt{--commit} ghi commit mã BFS, còn \texttt{--output} chỉ đến thư mục mới.
Manifest chứa SHA-256 đầu vào, trình đo, nhị phân và dữ liệu từng đồ thị.

\section{Tạo lại bảng, hình và báo cáo}
Quy trình cần Python với Matplotlib và \texttt{psutil}, cùng XeLaTeX và BibTeX.
Từ thư mục gốc, chạy lệnh sau để tái tạo các bảng, hình và PDF từ dữ liệu đã lưu.
\begin{quote}\small
\begin{verbatim}
make tables figures supplement-artifacts report
\end{verbatim}
\end{quote}
Khi bảng và hình đã có sẵn, có thể chạy riêng \texttt{make report}.
Kịch bản \path{scripts/benchmark_extended.py} dùng CLI trong tiến trình mới cho mỗi lần đo.
Giao thức này khác hai phiên trong báo cáo.

\chapter{Thống kê đầy đủ của từng cấu hình}
\label{app:quartiles}
Bảng~\ref{tab:quartiles} ghi thời gian của 18 cặp đồ thị và phiên bản.
Mỗi hàng dùng đúng 12 mẫu chính thức. Hai cột cực trị cho biết phạm vi mẫu,
còn cột IQR đo bề rộng của khoảng giữa hai tứ phân vị.
Các thống kê có thêm chữ số so với Bảng~\ref{tab:results} để thuận tiện đối chiếu.

\begin{table}[!htbp]
\centering\footnotesize
\setlength{\tabcolsep}{3pt}
\caption[Tứ phân vị và cực trị thời gian]{Tứ phân vị, độ rộng IQR và cực trị thời gian (ms)}
\label{tab:quartiles}
\begin{tabular}{@{}ll *{6}{S[table-format=4.5]}@{}}
\toprule
Đồ thị & Bản BFS & {$Q_1$} & {Trung vị} & {$Q_3$} & {IQR} & {Min} & {Max} \\
\midrule
% Generated by scripts/analysis/summarize.py
Chu trình & Serial & 0.01000 & 0.01045 & 0.01063 & 0.00063 & 0.00960 & 0.01160 \\
 & Basic & 293.74090 & 321.52060 & 370.42275 & 76.68185 & 270.97300 & 551.66900 \\
 & Optimized & 496.35743 & 516.12445 & 652.13908 & 155.78165 & 455.01090 & 918.03240 \\
Thang & Serial & 0.01538 & 0.01615 & 0.01695 & 0.00158 & 0.01320 & 0.01830 \\
 & Basic & 133.90773 & 136.30730 & 153.46785 & 19.56013 & 130.24470 & 196.04300 \\
 & Optimized & 223.70240 & 231.44605 & 250.62933 & 26.92693 & 216.28490 & 354.40120 \\
Lưới & Serial & 0.03380 & 0.03570 & 0.03625 & 0.00245 & 0.03250 & 0.06260 \\
 & Basic & 20.72235 & 21.57335 & 22.27610 & 1.55375 & 20.62440 & 26.44760 \\
 & Optimized & 34.12415 & 36.04625 & 37.19375 & 3.06960 & 33.79550 & 41.88110 \\
Cây bậc ba nhỏ & Serial & 0.08520 & 0.09515 & 0.10473 & 0.01953 & 0.07790 & 0.18370 \\
 & Basic & 2.69240 & 2.84440 & 3.03113 & 0.33873 & 2.56500 & 3.41070 \\
 & Optimized & 3.23913 & 3.28730 & 3.40553 & 0.16640 & 3.16710 & 3.68560 \\
Cây nhị phân lớn & Serial & 596.47653 & 1169.43130 & 1281.29180 & 684.81528 & 573.65780 & 1303.45390 \\
 & Basic & 704.05380 & 765.42725 & 805.26718 & 101.21338 & 679.38740 & 852.19750 \\
 & Optimized & 471.33528 & 496.79215 & 506.26693 & 34.93165 & 394.59400 & 609.66270 \\
Cây bậc bốn lớn & Serial & 161.29720 & 163.64135 & 169.03665 & 7.73945 & 154.82710 & 174.94160 \\
 & Basic & 213.69665 & 214.87830 & 217.28555 & 3.58890 & 210.74150 & 228.42610 \\
 & Optimized & 121.43543 & 130.36645 & 140.95940 & 19.52398 & 109.51530 & 198.17960 \\
\bottomrule

\end{tabular}
\end{table}

